% --- CORRECCIÓN IMPORTANTE EN LA PRIMERA LÍNEA ---
% Agregamos "xcolor=table" para que funcione \rowcolor sin errores
\documentclass[aspectratio=169, 10pt, xcolor=table]{beamer}

% --- Configuración del Tema (estilo profesional como el ejemplo) ---
\usetheme{Madrid}
\usecolortheme{dolphin}

% --- Paquetes necesarios ---
\usepackage[utf8]{inputenc}
\usepackage[spanish]{babel}
\usepackage{amsmath, amssymb}
\usepackage{booktabs}
\usepackage{graphicx}
\usepackage{listings}
\usepackage{tikz}
\usetikzlibrary{arrows.meta, positioning, shapes.geometric}

% --- Ruta de Imágenes (crea la carpeta "imagenes" y coloca ahí tus archivos) ---
\graphicspath{{./imagenes/}}

% --- Estilo para código Python ---
\definecolor{codegreen}{rgb}{0,0.6,0}
\definecolor{codegray}{rgb}{0.5,0.5,0.5}
\definecolor{codepurple}{rgb}{0.58,0,0.82}
\definecolor{backcolour}{rgb}{0.95,0.95,0.92}
\lstset{
    language=Python,
    backgroundcolor=\color{backcolour},
    commentstyle=\color{codegreen},
    keywordstyle=\color{blue},
    numberstyle=\tiny\color{codegray},
    stringstyle=\color{codepurple},
    basicstyle=\ttfamily\scriptsize,
    breaklines=true,
    frame=single,
    showstringspaces=false
}

% --- Pie de página personalizado ---
\makeatletter

\setbeamertemplate{footline}{
  \leavevmode%
  \hbox{%
  \begin{beamercolorbox}[wd=.333333\paperwidth,ht=2.25ex,dp=1ex,center]{author in head/foot}%
    \usebeamerfont{author in head/foot}\insertshortauthor
  \end{beamercolorbox}%
  \begin{beamercolorbox}[wd=.333333\paperwidth,ht=2.25ex,dp=1ex,center]{title in head/foot}%
    \usebeamerfont{title in head/foot}Sesión 7
  \end{beamercolorbox}%
  \begin{beamercolorbox}[wd=.333333\paperwidth,ht=2.25ex,dp=1ex,right]{date in head/foot}%
    \usebeamerfont{date in head/foot}Machine Learning \hfill \insertframenumber/\inserttotalframenumber \hspace*{0.3cm}
  \end{beamercolorbox}}%
  \vskip0pt%
}

\makeatother

% --- Datos de la portada ---
\title[SESIÓN 7]{\textbf{Sesión 7}}
\subtitle{Gradient Boosting: XGBoost, LightGBM y el estado del arte en tablas}
\author[Prof. C. Sánchez]{Carlos César Sánchez Coronel}
\institute{\textbf{MACHINE LEARNING}}
\date{}

% Logo opcional (descomentar si tienes la imagen)
% \titlegraphic{\includegraphics[height=1.5cm]{logo.png}}

% --- Eliminar la barra de navegación (opcional) ---
\setbeamertemplate{navigation symbols}{}

% --- Índice al inicio de cada sección ---
\AtBeginSection[]{
  \begin{frame}
    \frametitle{Contenido}
    \tableofcontents[currentsection]
  \end{frame}
}

% ============================================================================
% INICIO DEL DOCUMENTO
% ============================================================================
\begin{document}

% --- Portada ---
\begin{frame}
    \titlepage
\end{frame}

% --- Índice general ---
\begin{frame}{Contenido}
    \tableofcontents
\end{frame}

% ============================================================================
% SECCIÓN 1: INTRODUCCIÓN Y CONCEPTO DE BOOSTING
% ============================================================================
\section{Introducción y Concepto de Boosting}

\begin{frame}{De bagging a boosting}
    \begin{itemize}
        \item \textbf{Random Forest (bagging)}: árboles en paralelo, promedian predicciones → reducen varianza.
        \item \textbf{Boosting}: árboles en secuencia, cada uno corrige errores del anterior → reduce sesgo.
        \item La corrección iterativa permite alcanzar mayor precisión, pero con mayor riesgo de sobreajuste.
    \end{itemize}
    % Basado en Pregunta 1
    % IMAGEN: Esquema comparativo bagging vs boosting
\end{frame}

\begin{frame}{AdaBoost: el precursor}
    \begin{block}{Idea principal}
        \begin{enumerate}
            \item Asignar pesos iguales a todas las observaciones.
            \item Entrenar un clasificador débil (ej. árbol de profundidad 1, ``stump'').
            \item Aumentar el peso de las observaciones mal clasificadas.
            \item Repetir y combinar por votación ponderada.
        \end{enumerate}
    \end{block}
    \begin{itemize}
        \item \textbf{Limitaciones}: sensible a ruido y outliers; no se generaliza fácilmente a otras funciones de pérdida.
    \end{itemize}
    % Basado en Pregunta 2
\end{frame}

% ============================================================================
% SECCIÓN 2: GRADIENT BOOSTING MATEMÁTICO
% ============================================================================
\section{Gradient Boosting Matemático}

\begin{frame}{Marco general}
    Objetivo: minimizar $\mathcal{L}(F) = \sum_{i=1}^n L(y_i, F(x_i))$ con $L$ diferenciable.
    \begin{block}{Construcción secuencial}
        \[
        F_m(x) = F_{m-1}(x) + \eta \, h_m(x)
        \]
        \begin{itemize}
            \item $F_{m-1}$: modelo hasta la iteración anterior.
            \item $h_m$: árbol añadido en la iteración $m$.
            \item $\eta$: \textbf{tasa de aprendizaje (learning rate)}.
        \end{itemize}
    \end{block}
    % Basado en Pregunta 3
\end{frame}

\begin{frame}{Ajuste de residuos (pérdida cuadrática)}
    Para $L(y,F) = \frac{1}{2}(y-F)^2$, el gradiente negativo es:
    \[
    r_{im} = -\left.\frac{\partial L}{\partial F}\right|_{F=F_{m-1}(x_i)} = y_i - F_{m-1}(x_i)
    \]
    \begin{itemize}
        \item Cada nuevo árbol $h_m$ se entrena para predecir los \textbf{residuos} del modelo actual.
        \item Luego se suma escalado por $\eta$.
    \end{itemize}
    % Basado en Pregunta 3
\end{frame}

\begin{frame}{Pérdida general: pseudo-residuos}
    Para una pérdida arbitraria, se usan los \textbf{gradientes negativos}:
    \[
    g_{im} = -\left.\frac{\partial L(y_i, F(x_i))}{\partial F(x_i)}\right|_{F=F_{m-1}}
    \]
    \begin{enumerate}
        \item Calcular $g_{im}$ para cada observación.
        \item Ajustar un árbol $h_m$ a los pares $(x_i, g_{im})$.
        \item Actualizar $F_m(x) = F_{m-1}(x) + \eta \, h_m(x)$.
    \end{enumerate}
    Esto equivale a un paso de gradiente descendente en el espacio de funciones.
\end{frame}

\begin{frame}{Algoritmo genérico de Gradient Boosting}
    \begin{enumerate}
        \item Inicializar $F_0(x) = \arg\min_\gamma \sum_i L(y_i, \gamma)$ (una constante).
        \item Para $m = 1$ hasta $M$:
            \begin{enumerate}
                \item Calcular pseudo-residuos $g_{im}$.
                \item Ajustar árbol $h_m$ a $(x_i, g_{im})$.
                \item Determinar paso $\rho_m$ (o usar $\eta$ fijo).
                \item Actualizar $F_m = F_{m-1} + \eta \, h_m$.
            \end{enumerate}
        \item Salida $F_M(x)$.
    \end{enumerate}
    % Basado en teoría
\end{frame}

% ============================================================================
% SECCIÓN 3: IMPLEMENTACIONES MODERNAS
% ============================================================================
\section{XGBoost, LightGBM, CatBoost}

\begin{frame}{XGBoost (eXtreme Gradient Boosting)}
    \begin{itemize}
        \item \textbf{Regularización}: añade penalización a la pérdida:
        \[
        \mathcal{L} = \sum_i L(y_i,\hat{y}_i) + \sum_k \Omega(f_k), \quad \Omega(f) = \gamma T + \frac{1}{2}\lambda\|w\|^2 + \alpha|w|
        \]
        ($T$: número de hojas, $w$: pesos de hojas)
        \item \textbf{Manejo nativo de valores faltantes}: aprende dirección óptima.
        \item \textbf{Early stopping} para evitar sobreajuste.
        \item \textbf{Paralelización} en la búsqueda de splits.
    \end{itemize}
    % Basado en Pregunta 5
\end{frame}

\begin{frame}{LightGBM (Light Gradient Boosting Machine)}
    \begin{itemize}
        \item \textbf{GOSS} (Gradient-based One-Side Sampling): retiene observaciones con gradientes grandes y muestrea las de gradiente pequeño.
        \item \textbf{EFB} (Exclusive Feature Bundling): agrupa características mutuamente excluyentes.
        \item \textbf{Crecimiento leaf-wise} (por hoja) en lugar de level-wise: más preciso pero puede sobreajustar.
        \item \textbf{Soporte nativo para categóricas}.
    \end{itemize}
    % Basado en Pregunta 6
\end{frame}

\begin{frame}{CatBoost (Categorical Boosting)}
    \begin{itemize}
        \item Especializado en variables categóricas.
        \item Target encoding con ordenamiento aleatorio para evitar leakage.
        \item Árboles simétricos que aceleran la inferencia.
        \item Ideal cuando hay muchas categóricas.
    \end{itemize}
\end{frame}

% ============================================================================
% SECCIÓN 4: COMPARACIÓN CON RANDOM FOREST
% ============================================================================
\section{Comparación con Random Forest}

\begin{frame}{Random Forest vs Gradient Boosting}
    \begin{table}
        \centering
        \begin{tabular}{l|c|c}
            \toprule
            \textbf{Característica} & \textbf{Random Forest} & \textbf{Gradient Boosting} \\
            \midrule
            Construcción & Paralela (bagging) & Secuencial (boosting) \\
            Objetivo & Reducir varianza & Reducir sesgo \\
            Sobreajuste & Robusto por naturaleza & Requiere regularización \\
            Velocidad & Rápido (paralelo) & Puede ser más lento \\
            Número de árboles & Menos necesario & Puede necesitar miles \\
            Datos ruidosos & Robusto & Sensible \\
            Rendimiento típico & Bueno & Excelente (con ajuste) \\
            \bottomrule
        \end{tabular}
    \end{table}
    % Basado en Pregunta 1 y 14
\end{frame}

% ============================================================================
% SECCIÓN 5: HIPERPARÁMETROS CLAVE
% ============================================================================
\section{Hiperparámetros Clave}

\begin{frame}{Hiperparámetros comunes en XGBoost y LightGBM}
    \begin{columns}[t]
        \column{0.5\textwidth}
        \textbf{XGBoost}
        \begin{itemize}
            \item \texttt{learning\_rate} ($\eta$)
            \item \texttt{max\_depth}
            \item \texttt{subsample}
            \item \texttt{colsample\_bytree}
            \item \texttt{lambda} (L2)
            \item \texttt{alpha} (L1)
            \item \texttt{gamma}
        \end{itemize}
        \column{0.5\textwidth}
        \textbf{LightGBM}
        \begin{itemize}
            \item \texttt{learning\_rate}
            \item \texttt{num\_leaves} (relacionado con profundidad)
            \item \texttt{feature\_fraction}
            \item \texttt{bagging\_fraction}
            \item \texttt{lambda\_l1}, \texttt{lambda\_l2}
            \item \texttt{min\_data\_in\_leaf}
        \end{itemize}
    \end{columns}
    % Basado en Preguntas 9 y 10
\end{frame}

\begin{frame}{Ajuste de hiperparámetros}
    \begin{itemize}
        \item Relación $\eta$ y número de árboles: $\eta$ pequeño requiere más árboles, pero mejora generalización.
        \item \textbf{Early stopping} sobre conjunto de validación para elegir el número óptimo de árboles.
        \item Búsqueda con \texttt{RandomizedSearchCV} o \texttt{GridSearchCV}.
        \item Para LightGBM, controlar \texttt{num\_leaves} para evitar sobreajuste.
    \end{itemize}
    % Basado en Pregunta 4 y 10
\end{frame}

% ============================================================================
% SECCIÓN 6: MÉTRICAS DE EVALUACIÓN
% ============================================================================
\section{Métricas de Evaluación}

\begin{frame}{Clasificación: log-loss y AUC}
    \begin{block}{Log-loss (entropía cruzada)}
        \[
        \text{LogLoss} = -\frac{1}{n}\sum_{i=1}^n \left[ y_i \log(\hat{p}_i) + (1-y_i)\log(1-\hat{p}_i) \right]
        \]
        Penaliza predicciones seguras pero erróneas. Ideal para probabilidades bien calibradas.
    \end{block}
    \begin{block}{AUC-ROC}
        Mide la capacidad de ranking, independiente del umbral. Útil en problemas desbalanceados.
    \end{block}
    % Basado en Pregunta 11
\end{frame}

\begin{frame}{Regresión: RMSE, MAE, RMSLE}
    \begin{columns}[t]
        \column{0.33\textwidth}
        \textbf{RMSE}
        \[
        \sqrt{\frac{1}{n}\sum (y_i-\hat{y}_i)^2}
        \]
        Sensible a outliers.
        \column{0.33\textwidth}
        \textbf{MAE}
        \[
        \frac{1}{n}\sum |y_i-\hat{y}_i|
        \]
        Robusto a outliers.
        \column{0.33\textwidth}
        \textbf{RMSLE}
        \[
        \sqrt{\frac{1}{n}\sum (\log(y_i+1)-\log(\hat{y}_i+1))^2}
        \]
        Penaliza más la infraestimación.
    \end{columns}
    % Basado en Pregunta 12
\end{frame}

% ============================================================================
% SECCIÓN 7: APLICACIONES PRÁCTICAS
% ============================================================================
\section{Aplicaciones Prácticas}

\begin{frame}{Predicción de clics (CTR) en publicidad online}
    \begin{itemize}
        \item Datos masivos, alta cardinalidad, necesidad de velocidad.
        \item \textbf{LightGBM} es ideal por su eficiencia y manejo de categóricas.
        \item Métrica principal: log-loss o AUC.
    \end{itemize}
    % Basado en Pregunta 15
\end{frame}

\begin{frame}{Competencias Kaggle}
    \begin{itemize}
        \item La mayoría de soluciones ganadoras en datos tabulares usan XGBoost, LightGBM o ensembles.
        \item Ejemplos: House Prices, Otto Group, Porto Seguro.
    \end{itemize}
\end{frame}

\begin{frame}{Modelado de riesgo crediticio}
    \begin{itemize}
        \item Se necesita interpretabilidad y buen rendimiento.
        \item XGBoost proporciona importancia de características (gain).
        \item Se pueden usar pesos para clases desbalanceadas.
    \end{itemize}
\end{frame}

% ============================================================================
% SECCIÓN 8: CASO INTEGRADO: PREDICCIÓN DE CHURN
% ============================================================================
\section{Caso Integrado: Predicción de Churn}

\begin{frame}{Contexto y datos}
    \begin{itemize}
        \item Empresa de telecomunicaciones, 100,000 clientes, 50 variables.
        \item 15\% de abandono (clase positiva).
        \item Variables: demográficas, uso, reclamaciones, facturación.
    \end{itemize}
    \begin{block}{Preprocesamiento}
        \begin{itemize}
            \item Nulos: XGBoost los maneja internamente (no imputar).
            \item Categóricas: en LightGBM usar \texttt{categorical\_feature}.
            \item División entrenamiento/validación/prueba estratificada.
        \end{itemize}
    \end{block}
    % Basado en Pregunta 16
\end{frame}

\begin{frame}{Modelado y ajuste}
    \begin{itemize}
        \item Probar XGBoost y LightGBM con búsqueda aleatoria de hiperparámetros.
        \item Usar \textbf{early stopping} con métrica AUC en validación.
        \item Resultados típicos: LightGBM suele ser ligeramente superior.
    \end{itemize}
    \begin{table}
        \centering
        \begin{tabular}{l|c|c|c}
            \toprule
            Modelo & AUC & Log-loss & F1 (clase 1) \\
            \midrule
            XGBoost & 0.832 & 0.352 & 0.612 \\
            LightGBM & 0.845 & 0.341 & 0.631 \\
            Random Forest & 0.801 & 0.382 & 0.582 \\
            \bottomrule
        \end{tabular}
    \end{table}
\end{frame}

\begin{frame}{Importancia de características}
    \begin{itemize}
        \item Variables más importantes (gain):
        \begin{enumerate}
            \item Número de reclamaciones últimos 6 meses.
            \item Factura media.
            \item Antigüedad del cliente.
            \item Consumo de datos último mes.
        \end{enumerate}
        \item \textbf{Acciones de negocio}: reducir reclamaciones, ofertas personalizadas, programas de fidelización.
    \end{itemize}
    % Basado en Pregunta 13
\end{frame}

% ============================================================================
% SECCIÓN 9: IMPLEMENTACIÓN EN PYTHON
% ============================================================================
\section{Implementación en Python}

\begin{frame}[fragile]{XGBoost}
    \begin{lstlisting}
import xgboost as xgb
from sklearn.model_selection import train_test_split
from sklearn.metrics import roc_auc_score

X_train, X_val, y_train, y_val = train_test_split(X, y, test_size=0.2)

dtrain = xgb.DMatrix(X_train, label=y_train)
dval = xgb.DMatrix(X_val, label=y_val)

params = {
    'objective': 'binary:logistic',
    'eval_metric': 'logloss',
    'max_depth': 5,
    'learning_rate': 0.05,
    'subsample': 0.8,
    'colsample_bytree': 0.8,
    'lambda': 1.0,
    'alpha': 0.1
}

model = xgb.train(params, dtrain, num_boost_round=1000,
                  evals=[(dval, 'validation')],
                  early_stopping_rounds=50, verbose_eval=50)

y_pred = model.predict(dval)
auc = roc_auc_score(y_val, y_pred)
    \end{lstlisting}
\end{frame}

\begin{frame}[fragile]{LightGBM}
    \begin{lstlisting}
import lightgbm as lgb
from sklearn.metrics import roc_auc_score

train_data = lgb.Dataset(X_train, label=y_train)
val_data = lgb.Dataset(X_val, label=y_val, reference=train_data)

params = {
    'objective': 'binary',
    'metric': 'auc',
    'boosting_type': 'gbdt',
    'num_leaves': 31,
    'learning_rate': 0.05,
    'feature_fraction': 0.8,
    'bagging_fraction': 0.8,
    'bagging_freq': 5,
    'lambda_l1': 0.1,
    'lambda_l2': 1.0
}

model = lgb.train(params, train_data, valid_sets=[val_data],
                  num_boost_round=1000,
                  callbacks=[lgb.early_stopping(50), lgb.log_evaluation(50)])

y_pred = model.predict(X_val)
auc = roc_auc_score(y_val, y_pred)
    \end{lstlisting}
\end{frame}

% ============================================================================
% SECCIÓN 10: CASO AVANZADO: PREDICCIÓN DE COMPRA EN E-COMMERCE (Pregunta 20)
% ============================================================================
\section{Caso Avanzado: Predicción de Compra (E-commerce)}

\begin{frame}{Contexto del problema}
    \begin{itemize}
        \item Plataforma de e-commerce (ej. Mercado Libre).
        \item Predecir si un usuario comprará después de hacer clic (conversión).
        \item Tasa de conversión baja (2\%).
        \item Datos masivos, muchas variables categóricas (usuario, producto, contexto).
    \end{itemize}
\end{frame}

\begin{frame}{Feature engineering específico}
    \begin{enumerate}
        \item \textbf{Usuario}: frecuencia de compras, tasa de conversión histórica, promedio de gasto, días desde última compra.
        \item \textbf{Producto}: precio relativo a categoría, popularidad (clics/ventas), antigüedad.
        \item \textbf{Contexto}: hora (codificación cíclica), día de semana, dispositivo, ubicación.
        \item \textbf{Interacciones}: usuario × categoría, precio × dispositivo, tiempo en página × precio.
    \end{enumerate}
\end{frame}

\begin{frame}{Manejo del desbalance y evaluación}
    \begin{itemize}
        \item Usar \texttt{scale\_pos\_weight} en XGBoost o \texttt{class\_weight} en LightGBM.
        \item Métricas: AUC, log-loss, precisión/recall/F1 (ajustando umbral).
        \item Curva Precision-Recall especialmente informativa.
        \item Early stopping con conjunto de validación temporal.
    \end{itemize}
\end{frame}

\begin{frame}{Interpretabilidad y monitoreo}
    \begin{itemize}
        \item Importancia de características (gain) para entender factores globales.
        \item SHAP para explicar predicciones individuales a equipos de negocio.
        \item Monitoreo en producción: asociar métricas técnicas a negocio (recall $\leftrightarrow$ ventas perdidas, precisión $\leftrightarrow$ clics desperdiciados).
    \end{itemize}
\end{frame}

% ============================================================================
% SECCIÓN 11: CONCLUSIÓN
% ============================================================================
\section{Conclusión}

\begin{frame}{Lo que hemos aprendido}
    \begin{exampleblock}{Conceptos clave}
        \begin{itemize}
            \item Gradient Boosting: construcción secuencial, pseudo-residuos, tasa de aprendizaje.
            \item XGBoost: regularización, manejo de nulos, early stopping.
            \item LightGBM: GOSS, EFB, crecimiento leaf-wise, soporte categóricas.
            \item Comparación con Random Forest.
            \item Hiperparámetros y ajuste.
            \item Métricas: log-loss, AUC, RMSE, MAE, RMSLE.
            \item Casos prácticos: CTR, churn, e-commerce.
        \end{itemize}
    \end{exampleblock}

\end{frame}

% --- Diapositiva final ---
\begin{frame}
    \centering
    \Huge \textbf{¡Gracias!}

\end{frame}

\end{document}